\documentclass[hyperref={pdfpagelabels=false},table]{beamer}

\input{lec_common}
\date{}
\begin{document}

\part{Unit 8a: GUI using matplotlib}

%\frame{\tableofcontents}


\begin{frame}[fragile]{Widgets}
\vfill
Widget = windows gadget  (Engl.) is the term for a little window which serves as a user interface to
a computer program.
\vfill
We use here the widget module of Matplotlib
\begin{python}
 import matplotlib.widgets as mpw
\end{python}
This lecture explains the basic by a couple of guiding examples. Make sure you choose
in Spyder the graphic mode "Automatic".
\end{frame}

\begin{frame}[fragile]{Widgets and Events}
There are different use scenarios:
\begin{enumerate}
\item Using matplotlib.widgets we will make sliders and buttons in a widget.
\item In any plot (without using widgets) you can handle mouse events such as "mouse
pressed". We will make such an example.
\end{enumerate}
\end{frame}
\begin{frame}[fragile]{Creating a Widget}
In the last plot we created a figure object and an axes object in one step by the \pyth!subplots!-command.
Now we create them separately,
\begin{enumerate}
\item Make a figure object \pyth!fig = figure()!
\item Make an axes object within the figure \pyth!ax = axes(...)!
\item Create a widget, e.g. a slider,  by using this axes object \pyth!sld = Slider(ax, ...)!
\item Define a \emph{callback function}  that processes the data obtained from the widget.
\end{enumerate}
\end{frame}

\begin{frame}[fragile]{Axes dimensions - the size of the widget}
The signature of the \pyth!axes!-command:
\begin{python}
ax = axes([left, bottom, width, height])
\end{python}
\vfill
\includegraphics[height=0.6\textheight]{axesRectangle.png}
\end{frame}
\begin{frame}[fragile]{Turning a Rectangle into a Sliderobject}
 \begin{python}
sld = Slider(ax, 'amp', 0., 5.,
             valinit=1, valfmt="%1.1f")
\end{python}
The value choosen by the slider is available through \pyth!sld.val!
\end{frame}
\begin{frame}[fragile]{Putting things together}
 \begin{python}
from matplotlib.widgets import Slider
fig = figure()
sld_ax = axes([0.2, 0.9, 0.6, 0.05])
plt_ax = axes([0.1, 0.15, 0.8, 0.7])
sld = Slider(sld_ax, 'values', 0., 5.,
             valinit=1,  valfmt="%1.1f")
anno=plt_ax.annotate('Here comes later a plot',
             (0.2,0.40), fontsize=24)
 \end{python}
\begin{center}
\includegraphics[height=0.5\textheight]{slider1.png}
\end{center}
\end{frame}
\begin{frame}[fragile]{The callback function}
The callback function processes the given value from the slider.
\begin{python}
def val_annotate(val):
    anno.set_text(f'The given value is {val}')
\end{python}
Note, the call back function has only one parameter. All other information
is passed as a global variable -- here the annotation object \pyth!anno!.\\
The call back function is exectued on every change of the slider position by
\begin{python}
sld.on_changed(val_annotate)
\end{python}
\end{frame}
\begin{frame}[fragile]{The complete demonstration}
\begin{python}
from matplotlib.widgets import Slider

fig = figure()
sld_ax = axes([0.2, 0.9, 0.6, 0.05])
plt_ax = axes([0.1, 0.15, 0.8, 0.7])
val_init=1.0
sld = Slider(sld_ax, 'values', 0., 5.,
             valinit=val_init,  valfmt="%1.1f")
anno=plt_ax.annotate(f'The given value is'
          f' {val_init:1.1f}', (0.2,0.40),
                              fontsize=24)
def val_annotate(val):
    anno.set_text(f'The given value is {val:1.1f}')
sld.on_changed(val_annotate)
\end{python}
\end{frame}
\begin{frame}[fragile]{... and the result}
\includegraphics[width=\textwidth]{slider2.png}
\end{frame}
\begin{frame}[fragile]{Example 2 - A dynamic plot}
\begin{python}
from matplotlib.pyplot import *
from matplotlib.widgets import Slider
fig = figure(figsize = (4,2))
sld_ax = axes([0.2, 0.9, 0.6, 0.05]) # axes for slider
ax = axes([0.1, 0.15, 0.8, 0.7])     # axes for the plot
sld = Slider(sld_ax, 'amp', 0., 5.)
x = linspace(-2*pi, 2*pi, 200)
ax.set_ylim(-5.5, 5.5)
lines, = ax.plot(x, sld.val*sin(x))

def update_amplitude(val):
    lines.set_ydata(val*sin(x))

sld.on_changed(update_amplitude)
\end{python}
\end{frame}
\begin{frame}[fragile]{... and the result}
\includegraphics[width=\textwidth]{slider3.png}
\end{frame}
\begin{frame}[fragile]{Example 3 - Two sliders}
For the damped Lissajous curve defined by

\begin{eqnarray*}
        x(t) &= \cos(at)e^{-dt} \\
        y(t) &= \sin(t)e^{-dt}
\end{eqnarray*}
with $0\leq t < 14\pi$.


Use one slider for the angular frequency $a$ and one for the damping factor $d$.
\end{frame}
\begin{frame}[fragile]{Example 3 - The code}
\begin{python}
fig = figure()
slda_ax = axes([0.05, 0.5, 0.3, 0.03]) # axes for a
sldd_ax = axes([0.05, 0.4, 0.3, 0.03]) # axes for d
ax = axes([0.5, 0.1, 0.4, 0.8])        # axes for plot

t = linspace(0, 14*pi, 400)
a0, d0 = 1., 0.05
slda = Slider(slda_ax, 'a', 0, 2,
              valinit=a0, valfmt="%1.2f")
sldd = Slider(sldd_ax, 'd', 0, 0.1,
              valinit=d0,  valfmt="%1.3f")

ax.set_xlim([-1, 1])
ax.set_ylim([-1, 1])
ax.set_aspect(1.0) # the aspect ratio should be 1 here
line, = ax.plot(cos(a0*t)*exp(-d0*t), sin(t)*exp(-d0*t))
\end{python}
\end{frame}
\begin{frame}[fragile]{Example 3 - The code (Cont.)}
And finally the common callback function for both sliders
\begin{python}
def update_curve(val):
    a = slda.val
    d = sldd.val
    line.set_xdata(cos(a*t)*exp(-d*t))
    line.set_ydata(sin(t)*exp(-d*t))

slda.on_changed(update_curve)
sldd.on_changed(update_curve)
\end{python}
\end{frame}
\begin{frame}[fragile]{... and the result}
\includegraphics[width=\textwidth]{slider4.png}
\end{frame}
\end{document}
